\documentclass[11pt]{article}
\usepackage[margin=1in]{geometry}
\usepackage{amsmath,amssymb}
\usepackage{graphicx}

\title{GPU-Accelerated Pareto Frontier Enumeration for the\\
       Multi-Objective Traveling Salesperson Problem\\
       via Dynamic Layer Cutset on MDDs}
\author{}
\date{}

\begin{document}
\maketitle

% ============================================================
\section{Methodology}
% ============================================================

We describe the CPU baseline algorithm and the GPU-accelerated variant
for computing the exact Pareto frontier of a multi-objective Traveling
Salesperson Problem (TSP) using Multi-valued Decision Diagrams (MDDs)
and the dynamic layer cutset method.

% ------------------------------------------------------------
\subsection{Problem Setting}

Given an MDD $\mathcal{M}$ with $L$ layers encoding all feasible TSP
tours, each arc carries a weight vector
$\mathbf{w} \in \mathbb{Z}^p$ ($p$ objectives).
A root-to-terminal path corresponds to a tour with objective vector
equal to the sum of its arc weights.
The goal is to enumerate the \emph{Pareto frontier}---the set of all
non-dominated objective vectors---over all paths.

% ------------------------------------------------------------
\subsection{CPU Baseline: Dynamic Layer Cutset}

The CPU algorithm proceeds in three phases.

\subsubsection{Phase 1: Bidirectional Expansion}

Two sets of partial Pareto frontiers are grown from both ends of the MDD:

\begin{itemize}
  \item \textbf{Top-down:} Starting from the root node (layer~0) with
        the singleton frontier $\{\mathbf{0}\}$, expand layer by layer.
        At layer~$\ell$, for each node~$v$ and each incoming arc
        $(u \to v)$ with weight~$\mathbf{w}$, every point $\mathbf{x}$
        in $u$'s frontier generates a candidate $\mathbf{x} + \mathbf{w}$.
        Candidates arriving at the same node are merged and
        dominance-filtered to retain only non-dominated points.

  \item \textbf{Bottom-up:} Symmetric expansion from the terminal node
        (layer~$L{-}1$) towards the root, using outgoing arcs in reverse.
\end{itemize}

A \emph{layer-value heuristic}
$V(\ell) = \sum_{v \in \text{layer}(\ell)} |\text{frontier}(v)| \cdot |\text{arcs}(v)|$
estimates the expansion cost of the next layer.
At each step, the direction with the smaller layer value is expanded.
The two fronts meet at a \emph{cutset layer}~$c$ chosen dynamically
to minimise the total work.

\subsubsection{Phase 2: Cutset Convolution}

At the cutset layer, each node~$v$ has a top-down frontier
$F_v^{\downarrow}$ and a bottom-up frontier $F_v^{\uparrow}$.
The complete Pareto frontier is contained in
\[
  \bigcup_{v} \bigl\{ \mathbf{a} + \mathbf{b}
    : \mathbf{a} \in F_v^{\downarrow},\;
      \mathbf{b} \in F_v^{\uparrow} \bigr\}.
\]
The CPU processes nodes one-by-one:\ for each node~$v$, it forms all
$|F_v^{\downarrow}| \times |F_v^{\uparrow}|$ Cartesian-product sums
and incrementally merges them into a running global frontier, checking
each new candidate against the current frontier for dominance.
The per-candidate cost is $O(|F|)$ where $|F|$ is the current frontier
size, giving total complexity
$O\bigl(\sum_v |F_v^{\downarrow}| \cdot |F_v^{\uparrow}| \cdot |F|\bigr)$.

\subsubsection{Phase 3: Result}

The global frontier after processing all cutset nodes is the exact
Pareto frontier.

% ------------------------------------------------------------
\subsection{GPU-Accelerated Top-Down Variant}

We first discuss a purely top-down GPU-accelerated enumeration algorithm.
The key idea is to keep Pareto frontier data on the GPU throughout the entire
process, minimising host--device transfers.

\subsubsection{Phase 1: GPU Single-Direction Expansion}

\paragraph{Data packing.}
All MDD layers are packed into flat CSR-style arrays on the host and
transferred to the GPU once at startup.
For each layer~$\ell$, we store the incoming-arc offsets, source-node indices, 
and arc weights.

\paragraph{Expansion kernels.}
Each layer expansion is a three-step GPU pipeline:

\begin{enumerate}
  \item \textbf{Candidate generation} (\texttt{expand\_candidates\_points\_kernel}):
        Each warp processes one edge, reading source-node frontier points
        and adding the arc weight vector.
        This is embarrassingly parallel---every candidate is independent.

  \item \textbf{Per-node dominance filtering:}
        Candidates arriving at the same destination node are compared pairwise to filter out dominated points. We developed two versions of this kernel:
        \begin{itemize}
            \item \textbf{v1 (Dense 2D Grid, \texttt{mark\_dominated\_by\_dst\_tiled\_kernel}):}
                  A 2D grid assigns one block-row per node, launching \texttt{<<<num\_nodes, max\_blocks>>>}. 
                  Each thread loads its candidate into registers, then iterates over shared-memory tiles of all candidates in the same node. 
                  However, highly unbalanced MDD layers result in massive block scheduling overhead, as many nodes with zero candidates still force the GPU to schedule empty blocks.
            \item \textbf{v2 (Load-Balanced 1D Grid, \texttt{mark\_dominated\_1d\_kernel}):}
                  To solve the imbalance, we compute \texttt{ceil\_div(node\_candidates, threads)} for each node and perform a parallel prefix sum to find the exact total number of required blocks.
                  The kernel is launched with a 1D grid of exactly this size, completely eliminating empty block scheduling.
                  Within the kernel, a custom binary search on the prefix-sum array maps each 1D block to its correct destination node.
                  This perfect load balancing drastically reduces CUDA launch overhead, cutting total expansion time by over 30\%.
        \end{itemize}

  \item \textbf{Stream compaction} (\texttt{scatter\_alive\_kernel}):
        An exclusive prefix sum on the alive flags gives output indices;
        surviving points are scattered into a compact output buffer.
\end{enumerate}

\paragraph{Parallelism advantage.}
On the CPU, expansion is sequential over nodes and arcs.
On the GPU, all edges in a layer are expanded simultaneously (step~1),
and dominance filtering for all nodes in the layer runs in parallel
(step~2, one block-row per node).
For the largest layers (24\,024 nodes, 168\,168 edges, 4M+ points),
this yields $>$100$\times$ speedup over the CPU expansion alone.

\subsubsection{Phase 2: Result}
The global frontier at the terminal node is returned directly from device memory.


% ------------------------------------------------------------
\subsection{GPU-Accelerated Coupled Variant}

Building upon the highly optimized GPU expansion kernels described above,
we also ported the dynamic layer cutset (coupled) method to the GPU.

\subsubsection{Phase 1: GPU Bidirectional Expansion}
Data packing now encompasses both incoming arcs (for top-down) and outgoing arcs (for bottom-up).
The layer expansion logic directly references the \emph{Expansion kernels} from the GPU Top-Down Variant.

\paragraph{Layer-value heuristic.}
The heuristic $V(\ell)$ is computed on the GPU using a
\texttt{thrust::reduce}, returning a single scalar to the host.
This avoids copying the entire frontier back for the selection decision.

\subsubsection{Phase 2: GPU Batched Convolution}

The cutset convolution is the computational bottleneck for the coupled algorithm.
The CPU processes nodes sequentially; the GPU version uses a
\emph{batched frontier-filter pipeline}:

\begin{enumerate}
  \item \textbf{Batching.}
        Cutset nodes are sorted by decreasing product size and grouped
        into batches of up to 2\,M Cartesian products each.

  \item \textbf{Cartesian product} (\texttt{cartesian\_product\_kernel}):
        One CUDA block per node computes all $|F_v^{\downarrow}| \times
        |F_v^{\uparrow}|$ sums.  Threads stride over the product space.

  \item \textbf{Frontier filter}
        (\texttt{mark\_dominated\_by\_frontier\_kernel}):
        Each candidate is checked against the \emph{existing} frontier
        using shared-memory tiling.
        Cost: $O(\text{batch\_size} \times |F|)$.
        This is the key efficiency gain: with $|F| \approx 5$--$6$K and
        batch size $\approx 2$M, each thread checks $\sim$6K frontier
        points loaded tile-by-tile into shared memory.
        Since frontier points are loaded cooperatively into shared memory
        and reused by all 128 threads in a block, the effective memory
        bandwidth is amplified $128\times$.

        This step typically eliminates $>$99.9\% of candidates
        (e.g., 2M $\to$ $\sim$200--2000 survivors).

  \item \textbf{Frontier update}
        (\texttt{mark\_frontier\_dominated\_kernel}):
        Existing frontier points dominated by new survivors are removed.
        Cost: $O(|F| \times |\text{survivors}|)$.
        Survivors are appended to the frontier.

  \item \textbf{Final global dominance}
        (\texttt{mark\_dominated\_global\_kernel}):
        After all batches, a single $O(|F|^2)$ pass on the accumulated
        frontier ($\sim$10K points) cleans up any cross-batch dominated
        points.  This takes $<$2\,ms on the GPU.
\end{enumerate}

\paragraph{Parallelism advantage.}
The CPU's serial \texttt{convolute} checks each of the 303M candidate
products against the frontier one at a time.
The GPU checks 2M candidates \emph{simultaneously}, with each of the
$\lceil 2\text{M}/128 \rceil \approx 15$K thread blocks independently
scanning the $\sim$5K frontier via shared memory.
The frontier filter is thus the GPU analogue of the CPU's incremental
merge, but with $\sim$15K-way parallelism.

% ============================================================
\section{Results}
% ============================================================

We compare the CPU and GPU variants on two TSP instances
with 15 cities, using the dynamic layer cutset method (method~3).
All experiments were run on the same machine with a single NVIDIA GPU.

% ------------------------------------------------------------
\subsection{Summary}

\begin{table}[h]
\centering
\begin{tabular}{|l|r|r|r|r|r|}
\hline
\textbf{Instance} & \textbf{Objs} & \textbf{CPU Sols} & \textbf{CPU Time (s)} & \textbf{CUDA Time (s)} & \textbf{Speedup} \\
\hline
seed108   & 3 &  831 &  2.78 & 5.37 & 0.52$\times$ \\
seed11417 & 4 & 5372 & 18.24 & 5.94 & 3.07$\times$ \\
\hline
\end{tabular}
\caption{CPU vs.\ CUDA execution times on 15-city TSP instances.
         Solution counts match exactly in both cases.}
\label{tab:results}
\end{table}

% ------------------------------------------------------------
\subsection{Phase-Level Breakdown (4 Objectives)}

\begin{table}[h]
\centering
\begin{tabular}{|l|r|r|l|}
\hline
\textbf{Phase} & \textbf{CPU (s)} & \textbf{CUDA (s)} & \textbf{Notes} \\
\hline
MDD Packing     & ---   & 0.14  & One-time host$\to$device transfer \\
Expansion (TD)  & $\sim$0.7
                        & 0.12  & 9 layers, up to 4M points \\
Expansion (BU)  & (combined)      & 0.02  & 7 layers \\
Convolution     & 17.5  & 5.29  & 153 batches of $\sim$2M products \\
Final dominance & ---   & 0.001 & $O(10\text{K}^2)$ on GPU \\
\hline
\textbf{Total}  & 18.24 & 5.94  & \\
\hline
\end{tabular}
\caption{Phase breakdown for the 4-objective, 15-city instance.
         The convolution phase accounts for 96\% of CPU time
         but only 89\% of CUDA time.}
\label{tab:breakdown}
\end{table}

The GPU expansion phase is roughly $6\times$ faster than CPU for the
same work (0.14\,s vs.\ $\sim$0.7\,s), and the convolution phase
achieves a $3.3\times$ speedup (5.29\,s vs.\ 17.5\,s).

% ------------------------------------------------------------
\subsection{Why the 3-Objective Case is Slower on GPU}

For the 3-objective instance, the GPU variant is $\sim$2$\times$ slower
than the CPU.  Three factors explain this:

\begin{enumerate}
  \item \textbf{Small problem size.}
        The 3-objective instance has far fewer Pareto-optimal points at
        each MDD node (the Pareto front size scales super-polynomially
        in the number of objectives).
        The CPU finishes expansion and convolution in 2.78\,s total.
        Most of this time is in the convolution, which on the CPU involves
        only moderate-size frontiers ($\sim$800 solutions).

  \item \textbf{Fixed GPU overhead.}
        MDD packing (host$\to$device transfer of all 17 layers) and
        initial CUDA context setup take $\sim$0.8\,s regardless of
        problem difficulty.
        For a 2.78\,s CPU run, this overhead alone consumes 29\% of
        the CPU budget.

  \item \textbf{Underutilised parallelism.}
        With only $\sim$800 final solutions and small intermediate
        frontiers, the GPU batched convolution cannot amortise kernel
        launch overhead.
        The frontier filter kernel processes 2M candidates against a
        frontier of $\sim$800 points, but 15K thread blocks each scanning
        only $\lceil 800/128 \rceil = 7$ shared-memory tiles finish
        very quickly, making kernel launch and synchronisation overhead
        a large fraction of wall time.
        In contrast, for the 4-objective case the frontier is
        $\sim$5--6K, requiring $\sim$47 tiles per block, which keeps
        the GPU's streaming multiprocessors occupied.
\end{enumerate}

In summary, the GPU variant is designed for the regime where the
convolution dominates wall time---which occurs with 4+ objectives,
larger cities, or denser MDD structures.
For small instances where the CPU finishes in a few seconds, the fixed
overheads of GPU initialisation and data transfer outweigh the
parallelism gains.

\end{document}
